\section{Introduction}

\subsection{Motivation}

Producing a character-consistent animated episode with generative models
requires sustaining visual identity across dozens of shots, each one
regenerated from scratch by text-to-image and text-to-video models. In our
production of \nts, the observed one-pass acceptance rate was
\textbf{12 of 52 shots}
(\textbf{23\%}). The remaining 40 shots required re-generation, and a
non-trivial fraction were rejected more than once.

This figure is, to the best of our knowledge, \emph{not} unusual. It is,
however, almost never reported. Production teams treat per-shot failure as an
operational cost and have no incentive to publish its distribution. The
result is a field operating without a shared quantitative account of
\emph{what actually breaks}.

\subsection{Contributions}

We contribute a system description, an empirical failure taxonomy, and one
technique-level finding, packaged as a single industrial-grade case study.

\begin{itemize}
  \item \textbf{C1 -- System.}
        We describe \agnes, a shooting-stage pipeline that composes
        \comfy\ graph execution, a locked style specification
        (\styledef), and face-embedding-based character anchoring, and
        the standards that make its outputs reproducible across operators.

  \item \textbf{C2 -- Failure Analysis at episode scale.}
        We release the audit records of \nts: \textbf{52 shot-level
        inspection records} covering six scenes and 18 fields each, of
        which 40 were rejected. We provide the first public distribution
        of rejection reasons in an AIGC animation pipeline, together with
        a taxonomy derived from them.

  \item \textbf{C3 -- Negative-prompt efficacy.}
        We quantify a counter-intuitive behaviour: constraint expressed
        as a \emph{negative} prompt is systematically \emph{less} likely
        to suppress the artefact it targets than the same constraint
        expressed positively. We report this as an ablation with a
        single manipulated variable.
\end{itemize}

\subsection{Why the failure distribution matters}

Two consequences follow directly from a 23\% one-pass rate. First, per-shot
pricing models systematically overstate deliverable output. Second, automatic
quality control cannot be designed from the success cases alone: using a
motion metric as an automatic reject rule, we observe that
\textbf{36 of 52} records are classified as \emph{motion-anomalous} — a
classification with no discriminative power. The dominant defects lie
elsewhere, and identifying them is the substance of \Cref{sec:failure}.

\subsection{Paper structure}

\Cref{sec:related} reviews prior work; \Cref{sec:system} describes the
pipeline; \Cref{sec:failure} presents the failure taxonomy;
\Cref{sec:ablation} reports the negative-prompt ablation;
\Cref{sec:discussion} discusses implications; \Cref{sec:conclusion}
concludes.

% TODO(写作顺序建议)
%   1. 先写本节 Motivation 段（已含 23% 这个数字，是全篇最有力的一句）
%   2. 再写 03_system（素材最全，写起来最顺）
%   3. 再写 04_failure（数据已实跑，见下方 FIXME）
%   4. 05_ablation 必须等实验跑完才能写完整
%   5. 02_related_work 最后写（它决定 C3 是否还成立）
